\chapter{Einleitung}

\todo{Einleitung ausführen}

Kurze Einleitung, Einführung und Motivation des Problems

FPGAs werden in der Industrie eingesetzt, wo Geschwindigkeit und Energieeffizienz große Rolle spielen --> Beschleunigung der Prozesse und Kosten sparen ist sinnvoll --> FPGAs werden weiterentwickelt
Hardware teuer, Simulation günstig --> Wann braucht man Hardware, wann reicht auch Simulation? --> Möglichkeiten und Limitationen --> Performanceunterschiede zwischen echter Hardware und Simulation
Bekannt, dass Simulation langsamer ist, jedoch Größenordnung??

% Einige allgemeine Tipps zum Schreiben einer schriftlichen Ausarbeitung gibt es unter \url{https://hsp.pages.cs.uni-duesseldorf.de/programmierung/website/lectures/sopro/ausarbeitung/}

%Ein Kapitel sollte, wenn es Unterüberschriften enthält, immer mit einem einleitenden Text anfangen. Dieser Text sollte für die Leserschaft einordnen, warum das Kapitel in der gesamten Arbeit relevant ist. Warum sollte es gelesen werden?

%Bei manch einer Einleitung (insbesondere bei kürzeren Arbeiten) kann es gut sein, dass die Motivation schnell erklärt ist. Dann sollten die Überschriften 1.1 und 1.2 entfallen und die Einleitung kann direkt mit der Motivation starten.

%\section{Motivation}

%Das hier ist die Einleitung.
%Hier sollte darauf eingegangen werden, was die Motivation der Arbeit ist.
%Was ist der Kontext? Warum ist das, was du getan hast, interessant? Warum sollte ich das lesen?


\section{Aufbau der Arbeit}

Hintergrundinformationen zur Arbeit: genutzte Technologien, die man zum Verständnis der Arbeit benötigt (Hardware, Software, Tools)
Implementation: Beschreibung, was im technischen Teil der Bachelorarbeit gemacht wurde, der Prozess bis zu einem Setup, bei dem eine Performanceanalyse stattfinden konnte
Evaluation: Welche Benchmarks und welche Messungen genutzt wurden, was dadurch herausgefunden wurde und wie das ganze zu bewerten/einzuordnen ist
Andere Arbeiten, die sich mit einem ähnlichen Thema befasst haben
Zusammenfassung der Ergebnisse

%Gehe kurz darauf an, was den Leser in den nächsten Abschnitten erwartet.
%Die Gliederung kann vorab mit der betreuenden Person abgesprochen werden.
